\documentclass[12pt, a4paper]{article}
\usepackage[utf8]{inputenc}
\usepackage{amsmath, amssymb, amsthm, amsfonts,amsthm,tikz-cd}
\usepackage{marvosym}
\usepackage{mathrsfs}
\usepackage{CJKutf8}
\usepackage{bm}
\usepackage[margin=1in]{geometry}
\newcommand{\id}{\operatorname{id}}
\newcommand{\qz}{\mathbb{Q}/\mathbb{Z}}
\newcommand{\z}{\mathbb{Z}}
\renewcommand{\hom}{\operatorname{Hom}}
\theoremstyle{remark}
\newtheorem*{remark}{Remark}
\newenvironment{lemma}[1]{
    \noindent \textbf{Lemma #1.} \itshape
}{
    \vspace{1em}
}

% 重新定义 solution 环境，保留标识并在末尾留出 8cm 空白
\newenvironment{solution}
  {\paragraph{\textit{Solution:}}\par\medskip}
  {\hfill$\blacksquare$\vspace{0.1cm}}

\title{Abstract Algebra III: Assignment - Week 13}
\author{
    \begin{CJK*}{UTF8}{gbsn}
    钟皓宇 (12533556)
    \end{CJK*}
}
\date{\today}

\begin{document}

\maketitle
\section*{Problem 1}

Write out a proof of Maschke's Theorem in the case of representations over
\(\mathbb{C}\) along the following lines.
\begin{enumerate}
    \item Given a representation
\[
\rho:G\to GL(V),
\]
where \(V\) is a vector space over \(\mathbb{C}\), let \((\cdot,\cdot)\) be any
positive definite Hermitian form on \(V\). Define a new form
\((\cdot,\cdot)_1\) on \(V\) by
\[
(v,w)_1=\frac{1}{|G|}\sum_{g\in G}(gv,gw).
\]
Show that \((\cdot,\cdot)_1\) is a positive definite Hermitian form, preserved
under the action of \(G\); that is,
\[
(v,w)_1=(gv,gw)_1
\]
always.
    \item If \(W\) is a subrepresentation of \(V\), show that
\[
V=W\oplus W^\perp
\]
as representations.
\end{enumerate}
\begin{solution}
(1) For \(a,b\in \mathbb C\),
\[
(av_1+bv_2,w)_1
=\frac{1}{|G|}\sum_{g\in G}(a\,gv_1+b\,gv_2,gw)
=a(v_1,w)_1+b(v_2,w)_1,
\]
and similarly the second variable is conjugate-linear. Moreover,
\[
(w,v)_1
=\frac{1}{|G|}\sum_{g\in G}(gw,gv)
=\overline{\frac{1}{|G|}\sum_{g\in G}(gv,gw)}
=\overline{(v,w)_1}.
\]
Thus \((\cdot,\cdot)_1\) is Hermitian.
Next, for \(v\neq 0\), since each \(g\in G\) acts invertibly, we have \(gv\neq 0\). Hence
\[
(v,v)_1=\frac{1}{|G|}\sum_{g\in G}(gv,gv)>0.
\]
Therefore \((\cdot,\cdot)_1\) is positive definite.
Finally, we show that this form is preserved by the action of \(G\). Let \(h\in G\). Then
\[
(hv,hw)_1
=\frac{1}{|G|}\sum_{g\in G}(ghv,ghw).
\]
As \(g\) runs over \(G\), so does \(gh\). Therefore
\[
(hv,hw)_1
=\frac{1}{|G|}\sum_{u\in G}(uv,uw)
=(v,w)_1.
\]
Thus
\[
(v,w)_1=(hv,hw)_1
\]
for all \(h\in G\), so \((\cdot,\cdot)_1\) is \(G\)-invariant.

(2) Choose a basis of \(V\), so that \(V\simeq \mathbb C^n\), and take the
standard positive definite Hermitian form
\[
(v,w)=\sum_i v_i\overline{w_i}.
\]
Averaging this form over the \(G\)-action gives a new Hermitian form
\[
(v,w)_1=\frac{1}{|G|}\sum_{g\in G}(gv,gw).
\]
By the previous part, \((\cdot,\cdot)_1\) is positive definite and
\(G\)-invariant.

Now let \(W\subset V\) be a subrepresentation, and define
\[
W^\perp=\{v\in V\mid (v,w)_1=0 \text{ for all } w\in W\}.
\]
Since \((\cdot,\cdot)_1\) is positive definite, we have
\[
V=W\oplus W^\perp
\]
as vector spaces.

It remains to show that \(W^\perp\) is \(G\)-stable. Let \(u\in W^\perp\),
\(h\in G\), and \(w\in W\). Since \(W\) is a subrepresentation,
\(h^{-1}w\in W\). Using \(G\)-invariance of \((\cdot,\cdot)_1\), we get
\[
(hu,w)_1=(u,h^{-1}w)_1=0.
\]
Hence \(hu\in W^\perp\), so \(W^\perp\) is a subrepresentation. Therefore
\[
V=W\oplus W^\perp
\]
as representations.
\end{solution}




\section*{Problem 2}

Prove the following theorem of Burnside: let \(G\) be a finite group, \(k\) an
algebraically closed field in which \(|G|\) is invertible, and
\[
\rho:G\to GL(V)
\]
be a representation over \(k\). By taking a basis of \(V\), write each
endomorphism \(\rho(g)\) as a matrix. Let \(\dim V=n\). Show that the
representation is simple if and only if there exist \(n^2\) elements
\(g_1,\dots,g_{n^2}\) of \(G\) so that the matrices
\[
\rho(g_1),\dots,\rho(g_{n^2})
\]
are linearly independent, and that this happens if and only if the algebra
homomorphism
\[
kG\to \operatorname{End}_k(V)
\]
is surjective. Note that \(\rho\) itself is generally not surjective.


\begin{solution}
    Let
\[
A=\operatorname{span}_k\{\rho(g):g\in G\}\subseteq \operatorname{End}_k(V).
\]
Then \(A\) is the image of the algebra homomorphism
$
kG\to \operatorname{End}_k(V).
$


Suppose that there exist \(g_1,\dots,g_{n^2}\in G\) such that
\(\rho(g_1),\dots,\rho(g_{n^2})\) are linearly independent. Since
\(\dim_k\operatorname{End}_k(V)=n^2\), they form a basis of
\(\operatorname{End}_k(V)\). Hence every elementary matrix \(E_{ij}\), sending
\(e_j\) to \(e_i\) and all other basis vectors to \(0\), can be written as a
\(k\)-linear combination of the matrices \(\rho(g)\).
Let \(0\neq W\subseteq V\) be a subrepresentation and choose
\(0\neq v=\sum_{\ell=1}^n c_\ell e_\ell\in W\). Pick \(j\) with \(c_j\neq 0\).
Then for every \(i\),
\[
\frac{1}{c_j}E_{ij}v=e_i.
\]
Since \(W\) is stable under every \(\rho(g)\), it is stable under their
linear combinations. Thus \(e_i\in W\) for all \(i\), so \(W=V\). Therefore
\(V\) is simple.


Conversely, suppose that \(V\) is simple. Let
\[
A=\operatorname{im}(kG\to \operatorname{End}_k(V)).
\]
Since \(|G|\) is invertible in \(k\), Maschke's theorem implies that \(kG\) is
semisimple. Hence its quotient \(A\) is also semisimple.

Since \(V\) is simple as a \(G\)-representation, it is simple as an
\(A\)-module. By Schur's lemma, because \(k\) is algebraically closed,
\[
\operatorname{End}_A(V)=k.
\]
By the Artin--Wedderburn theorem, \(A\) is a product of matrix algebras over
\(k\). Since \(V\) is simple, only one simple factor acts nontrivially on
\(V\). But \(A\) is defined as the image in \(\operatorname{End}_k(V)\), so its
action on \(V\) is faithful. Therefore \(A\) has only this one factor, and its
action on \(V\) identifies it with
\[
\operatorname{End}_k(V).
\]
Thus
\[
A=\operatorname{End}_k(V).
\]
Hence \(\dim_k A=n^2\). Since \(A\) is spanned by the matrices \(\rho(g)\), one
can choose \(n^2\) linearly independent matrices among them.
\end{solution}



\section*{Problem 3}

Let
\[
G=\langle x\mid x^n=1\rangle
\]
be a cyclic group of order \(n\), and let \(\zeta_n\in \mathbb{C}\) be a
primitive \(n\)-th root of unity. Then the simple complex characters of \(G\)
are the \(n\) functions
\[
\chi_r(x^s)=\zeta_n^{rs},
\]
where \(0\le r\le n-1\).
\begin{solution}
    Let \(V\) be a simple complex representation of \(G\), and set
\(T=\rho(x)\). Since \(x^n=1\), we have \(T^n=I\). Hence the minimal
polynomial of \(T\) divides \(t^n-1\). Over \(\mathbb C\), the polynomial
\(t^n-1=\prod_{r=0}^{n-1}(t-\zeta_n^r)\) has distinct roots, so \(T\) is
diagonalizable and its eigenvalues are among
\(\zeta_n^0,\dots,\zeta_n^{n-1}\).

For an eigenvalue \(\lambda\), the eigenspace
\(V_\lambda=\{v\in V\mid Tv=\lambda v\}\) is \(G\)-stable, because
\(G=\langle x\rangle\). Since \(V\) is simple, only one eigenspace can be
nonzero. Thus \(V=V_\lambda\) for some \(\lambda\), so \(\dim V=1\).

Therefore every simple complex representation of \(G\) is one-dimensional.
Since \(\rho(x)^n=1\), we have \(\rho(x)=\zeta_n^r\) for some
\(0\le r\le n-1\). Hence, for every \(s\),
\[
\chi_r(x^s)=\rho(x^s)=\rho(x)^s=(\zeta_n^r)^s=\zeta_n^{rs}.
\]
These \(n\) characters are distinct since
\(\chi_r(x)=\zeta_n^r\). Hence they are exactly all the simple complex
characters of \(G\).
\end{solution}
\end{document}